\documentclass[10pt]{article}
\usepackage[margin=1in]{geometry}
\usepackage{longtable,booktabs,array}
\usepackage{microtype}
\usepackage[hidelinks]{hyperref}
\usepackage{titlesec}
\titleformat{\section}{\large\bfseries}{}{0pt}{}
\setlength{\parindent}{0pt}\setlength{\parskip}{4pt}
\renewcommand{\arraystretch}{1.15}
\title{\textbf{PolyStep for Operations Research: Faithful-Reuse Report}}
\author{}\date{}
\begin{document}\maketitle\vspace{-2em}

\section*{Faithful-Reuse Audit: PolyStep Decision-Focused-Learning Study}

Project root: \texttt{/media/anindex/Data/ot-or-project}. This report documents, for each experiment, which upstream or
author code was reused and from where, exactly which parts were changed locally and why, and whether the run used
the authors' paper-default configuration. The governing rule audited is: for any benchmark whose authors released
code, reuse THEIR code at the PAPER-DEFAULT setting (cross-checked against code defaults); implement from scratch
ONLY where no author code exists; and bring the PolyStep method ONLY from its own pristine repository.

\section*{Top-level summary}

\begin{itemize}\setlength{\itemsep}{2pt}
\item \textbf{PolyStep is used unmodified from its own repository.} The library at \texttt{/media/anindex/Data/ot-or-project/polystep}
\end{itemize}
  has an empty working tree (\texttt{git -C polystep status --porcelain} returns nothing), tracks upstream
  \texttt{https://github.com/anindex/polystep.git}, and sits at commit \texttt{3c1d14f}. Every Python experiment imports it the
  same way, \texttt{sys.path.insert(0, "polystep/src"); from polystep import PolyStepOptimizer} (plus
  \texttt{polystep.epsilon.CosineEpsilon}, \texttt{polystep.hybrid\_subspace}, \texttt{polystep.transform}). One script,
  \texttt{exp\_spotree\_polystep.py}, uses the lower-level \texttt{from polystep.solver import PolyStep} entry point for a
  free-function objective, still from the same pristine package. No local copy or edit of PolyStep math exists.

\begin{itemize}\setlength{\itemsep}{2pt}
\item \textbf{Benchmarks that reuse authors' code verbatim (reused-as-is):}
\end{itemize}
  the entire PyEPO library (data generators for shortest-path / knapsack / TSP / portfolio, the Gurobi optimization
  models, the \texttt{pyepo.metric.regret} evaluator, and the full surrogate/DFL baseline suite SPO+, DBB, DPO, IMLE, PFYL,
  NCE, LTR, PG); the Globe-TSP Gurobi cutting-plane solver and DBB autograd function from \texttt{martius-lab/blackbox-backprop};
  the Blossom V min-cost-matching solver from \texttt{blossom\_python}; the ICON benchmark (data, \texttt{SolveICON} Gurobi ILP, and
  all three baselines from \texttt{predopt-benchmarks}); the IEEE-CIS MA\&RE Gurobi MILP optimiser and data loaders; the Amazon
  ALMRRC OR-tools reference router and official \texttt{rc-cli} scorer; the SPOTree predictor and its Gurobi shortest-path
  solver; and the ODECE knapsack data generators (generator only, see below).

\begin{itemize}\setlength{\itemsep}{2pt}
\item \textbf{Benchmarks adapted from authors' artifacts (adapted):}
\end{itemize}
  the DistrictNet Julia codebase (upstream code reused, with a new PolyStep estimator, feasibility repair, and reduced
  compute budgets, detailed below); the Toni / DSIRP inventory-routing column (authors' released instances and trained
  GLM weights reused, with a uniform-demand branch added); the ODECE MDKP experiments (authors' generator reused, their
  solver and loss NOT run, their published numbers cited from \texttt{Results/}); and the POLSKA max-flow and ABILENE
  min-cost-vertex-cover experiments (authors' shipped topology and instance files reused, objective/correction-function
  definitions ported, batched solvers reimplemented).

\begin{itemize}\setlength{\itemsep}{2pt}
\item \textbf{Benchmarks implemented from the paper because no reusable author code exists (original-following-paper):}
\end{itemize}
  SFGE (a score-function / REINFORCE estimator following Silvestri et al., JAIR 2024; no SFGE package exists); the
  synthetic prediction-in-constraint problems (uncertain-consumption knapsack, the L0 to L4 constraint ladder, the
  MDKP-consumption family, the WSMC set-multicover); the GPU-batched constraint solvers in \texttt{pto/solvers.py} (each
  validated against Gurobi or cvxpy); the Warcraft Bellman-Ford solver and CNN (reimplemented to the paper, validated
  against the dataset's stored optima); the IRP predict to CPCTSP to simulate pipeline (a faithful Python port of
  Greif et al.'s released Julia, plus an original GPU-batched exact CPCTSP solver validated against Gurobi); and all
  oracle-cost and scaling measurement harnesses.

A single methodological caveat applies to DistrictNet: the local runs reduced several compute budgets below the
paper defaults (Gurobi/ILS \texttt{MAX\_TIME} 1200 s to 120 s, GNN epochs cut, evaluation \texttt{NB\_SCENARIO} 10 to 50). These are
documented per file below so the report can state the deviation plainly.

\medskip\hrule\medskip

\section*{Shared harness provenance (\texttt{pto/})}

The Python experiments share the \texttt{pto/} package. Its problem and baseline provenance underpins most entries below.

\begin{itemize}\setlength{\itemsep}{2pt}
\item \textbf{Problems.} \texttt{pto/capability.py} (lines 14-18) imports the canonical four problems directly from PyEPO:
\end{itemize}
  \texttt{pyepo.data.\{shortestpath,knapsack,tsp,portfolio\}}, `pyepo.model.grb.\{shortestPathModel,knapsackModel,tspMTZModel,
  portfolioModel\}\texttt{, }pyepo.data.dataset.optDataset\texttt{, and }pyepo.func\texttt{. }pto/problems.py` adds a parallel set:
  \texttt{ShortestPath} (also PyEPO-backed) plus the original synthetic generators \texttt{TerrainSP}, \texttt{KnapsackWeights}, and
  \texttt{MDKPConsumption}.
\begin{itemize}\setlength{\itemsep}{2pt}
\item \textbf{Baselines.} SPO+ is PyEPO's native \texttt{pyepo.func.SPOPlus} (Elmachtoub and Grigas 2022); the full DFL suite (DBB,
\end{itemize}
  NID, DPO, IMLE, PFYL, NCE, LTR, PG) is taken from \texttt{pyepo.func}. None of these is reimplemented. \texttt{two-stage} is a
  trivial local MSE trainer. \texttt{SFGE} is an original score-function trainer following Silvestri et al. (JAIR 2024), with
  no author package available.
\begin{itemize}\setlength{\itemsep}{2pt}
\item \textbf{Solvers and accounting.} \texttt{pto/solvers.py} holds our batched GPU solvers (knapsack DP, fractional knapsack, MDKP
\end{itemize}
  greedy, set-multicover, max-flow, etc.), each validated against Gurobi or cvxpy. \texttt{pto/budget.py} provides the
  \texttt{SolveCounter}, \texttt{Timer}, and Gurobi-solve accounting used by the scaling and oracle-cost studies.

\medskip\hrule\medskip

\section*{Core benchmark / methods experiments}

These run on PyEPO synthetic data with PyEPO baselines, so they are classified as reused-as-is for the PyEPO problems
and baselines, original-following-paper for SFGE, and pristine import for PolyStep. All are warm-started from the same
two-stage model unless noted.

\begin{itemize}\setlength{\itemsep}{2pt}
\item \textbf{\texttt{challenge\_established.py}.} Data: PyEPO shortest-path / knapsack / TSP / portfolio via \texttt{pto.capability} setups.
\end{itemize}
  Solver: PyEPO Gurobi models for the regret metric, batched GPU solvers inside the PolyStep/SFGE closures. Methods:
  two-stage, SPO+/PFYL/IMLE (PyEPO), SFGE, PolyStep, all sharing one two-stage warm start (fairness protocol). Config:
  deg in \{2,4,6\}, seeds \{0,1,2\}, DFL warm epochs 30, PolyStep 150 steps, SFGE 120 epochs / 8 samples, batch 128. PyEPO
  data generators at their defaults.

\begin{itemize}\setlength{\itemsep}{2pt}
\item \textbf{\texttt{exp\_fair\_batched\_spo.py}.} Adds a second SPO+ backend: "SPO+" is PyEPO's per-instance Gurobi \texttt{F.SPOPlus}, while
\end{itemize}
  "SPO+(batched)" reimplements the exact SPO+ subgradient (\texttt{\_BatchedSPOPlusFunc}) on the shared batched GPU oracle,
  mirroring \texttt{pyepo.func.surrogate.SPOPlusFunc} term-for-term and checked equal to Gurobi by \texttt{verify\_oracle}. The batched
  variant is an adapted backend of PyEPO's algorithm, not a new method. Config FULL = \{SPO 30 epochs, PolyStep 150
  steps, SFGE 120 epochs / 8 samples\}, lr 1e-2, deg 4, seeds \{0,1,2\}.

\begin{itemize}\setlength{\itemsep}{2pt}
\item \textbf{\texttt{exp\_polystep\_vs\_sfge.py}.} PyEPO shortest-path and knapsack only. SFGE versus PolyStep head-to-head with
\end{itemize}
  two-stage / SPO+ / IMLE as references. Uses per-method tuned hyperparameters in \texttt{BEST} from prior LR sweeps and
  \texttt{pto.sweep\_lr.train\_spoplus\_lr}. Config: knapsack dims \{10,20,50,100,200\}, seeds \{0..4\}, plus sigma / probe-radius
  sensitivity grids. PyEPO data at defaults; PolyStep pristine.

\begin{itemize}\setlength{\itemsep}{2pt}
\item \textbf{\texttt{exp\_polystep\_bench.py}.} PyEPO four problems, three polytopes (orthoplex / simplex / cube). Compares the full
\end{itemize}
  PyEPO DFL suite plus two-stage plus PolyStep, with PolyStep hyperparameters loaded per (problem, polytope) from
  \texttt{exp\_results/polystep\_tune.json}. Wilcoxon(PolyStep < SPO+). 5 seeds, deg 4. PyEPO baselines reused-as-is.

\begin{itemize}\setlength{\itemsep}{2pt}
\item \textbf{\texttt{exp\_polystep\_tune.py}} and \textbf{\texttt{exp\_polystep\_ablation.py}.} PolyStep-only utilities over the PyEPO setups.
\end{itemize}
  \texttt{tune} grids probe count \{1,2,4\} and radii and writes the JSON the bench consumes; \texttt{ablation} sweeps polytope x probe
  over 5 seeds. Both warm-start from two-stage. Pristine PolyStep on PyEPO synthetic data.

\begin{itemize}\setlength{\itemsep}{2pt}
\item \textbf{\texttt{exp\_sfge\_suite.py}.} Two problems straight from the SFGE paper (Silvestri et al., JAIR 2024): unknown-weight
\end{itemize}
  knapsack and weighted set multi-cover, with prediction in the constraints. Features come from
  \texttt{pyepo.data.knapsack.genData}; the solvers are ours (\texttt{solve\_fractional\_knapsack} checked vs cvxpy, \texttt{solve\_set\_multicover}
  gap-checked vs a Gurobi exact ILP). SPO+/IMLE/PFYL/DBB declared not-applicable because the optimization-oracle camp is
  undefined for constraint prediction. Classification: original problem and solver implementation following the SFGE
  paper, original SFGE, pristine PolyStep.

\begin{itemize}\setlength{\itemsep}{2pt}
\item \textbf{\texttt{exp\_spo\_convergence.py}.} Epoch-grid convergence of SPO+ (PyEPO native) and SFGE, with a tuned PolyStep reference
\end{itemize}
  line, all warm-started from one two-stage model. Reviewer-inoculation sweep, not a new comparison. Epoch grid
  \{5..160\}, seeds \{0..4\}, deg 4.

\begin{itemize}\setlength{\itemsep}{2pt}
\item \textbf{\texttt{cross\_sfge.py}.} Two-stage, SPO+ (PyEPO via \texttt{train\_dfl}), SFGE, PolyStep on the four PyEPO problems, all
\end{itemize}
  warm-started from two-stage, capability defaults, seeds \{42,43,44\}.

\begin{itemize}\setlength{\itemsep}{2pt}
\item \textbf{\texttt{fair\_init.py}.} A fairness control on PyEPO knapsack only: compares cold (all from scratch) versus warm (all from
\end{itemize}
  two-stage) initialization across SPO+ (PyEPO), SFGE, PolyStep. Not a benchmark. Seeds \{0,1,2\}, deg 4, lr 1e-2.

\medskip\hrule\medskip

\section*{Prediction-in-constraint / harder experiments}

\begin{itemize}\setlength{\itemsep}{2pt}
\item \textbf{\texttt{blackbox\_constraint.py}.} Data: PyEPO \texttt{knapsack.genData} features, with the predicted quantity placed in the
\end{itemize}
  capacity constraint (a construction original to this project; no author code). Solver: ours, exact batched 0/1
  knapsack DP (\texttt{pto.solvers.knap1\_dp}). Methods: two-stage, SFGE, PolyStep (SPO+/PFYL/IMLE structurally not-applicable).
  Classification: original-following-paper on the PyEPO knapsack generator (used at defaults).

\begin{itemize}\setlength{\itemsep}{2pt}
\item \textbf{\texttt{exp\_constraint\_complexity.py}.} Data: a self-contained synthetic generator with an L0 to L4 constraint-type
\end{itemize}
  ladder (box, linear, cardinality, conic, prediction-coupled), the script's own design. Solver: ours, batched GPU
  greedy per level, verified against a Gurobi exact MIQCP. Methods: two-stage / SFGE / PolyStep with sigma and
  probe-radius fragility sweeps. Classification: original implementation (no released author code).

\begin{itemize}\setlength{\itemsep}{2pt}
\item \textbf{\texttt{exp\_constraint\_scaling.py}.} Two regimes: the constraint regime uses \texttt{pto.problems.MDKPConsumption} (original
\end{itemize}
  synthetic generator, solver \texttt{pto.solvers.mdkp\_greedy}); the contrast regime uses PyEPO \texttt{knapsack.genData} +
  \texttt{knapsackModel} so SPO+ applies. Classification: original for the MDKP-consumption regime, reused-as-is for the PyEPO
  SPO+ contrast.

\begin{itemize}\setlength{\itemsep}{2pt}
\item \textbf{\texttt{exp\_manyconstraints.py}.} Regime 1 is PyEPO grid shortest-path (\texttt{pto.problems.ShortestPath} + PyEPO Gurobi DAG
\end{itemize}
  solver + SPO+); regime 2 is the synthetic \texttt{MDKPConsumption}. Classification: adapted (PyEPO shortest-path reused) plus
  original (MDKP).

\begin{itemize}\setlength{\itemsep}{2pt}
\item \textbf{\texttt{exp\_odece\_full.py}} and \textbf{\texttt{exp\_odece\_mdkp.py}.} Data: the ODECE author generators are reused as-is,
\end{itemize}
  \texttt{from OptProblems.knapsack.kpdata import genCapacity, genWeights} (from \texttt{baselines/odece\_neurips25}), at settings
  matched to ODECE's shell scripts (\texttt{MDKP\_CapaExp.sh}, \texttt{MDKP\_WeightExp.py}): 50 items, dim 3, deg 6, noise 0.25, seeds
  \{11..15\}, capacity ratio \{0.18,0.20,0.22\}. Solver: OURS, \texttt{pto.solvers.mdkp\_greedy} + \texttt{mdkp\_repair} for deploy and a
  local Gurobi MDKP for the true-optimum oracle. ODECE's own solver and loss are NOT imported or run (run separately in
  ODECE's pinned environment); ODECE's published metrics are cited from \texttt{baselines/odece\_neurips25/Results/}. Methods:
  two-stage / SFGE / PolyStep live, ODECE and MSE numbers tabulated from \texttt{Results/}. Classification: adapted (problem
  generator reused as-is, solver and methods original). Paper-default settings, matching the ODECE scripts.

\begin{itemize}\setlength{\itemsep}{2pt}
\item \textbf{\texttt{exp\_maxflow\_polska.py}.} Data: the POLSKA topology is loaded from the shipped file
\end{itemize}
  \texttt{baselines/cpaior23\_branch\_and\_learn/Max Flow/data/graph\_POLSKA.txt} (12 nodes, 18 edges, not hardcoded) with per-fold
  \texttt{train/test\_POLSKA(k).txt} instances. Solver: ours, a batched GPU max-flow plus an exact batched port of the authors'
  Correction-Function-A, validated to <= 1e-3 against a numpy reference and Gurobi LP; the authors' C++ binaries are not
  invoked. Methods: two-stage / SFGE / PolyStep with a per-edge \texttt{Linear(8,1)} predictor mirroring \texttt{Ridge\_para\_POLSKA}.
  Classification: adapted (authors' data and correction-function reused faithfully, solver reimplemented).

\begin{itemize}\setlength{\itemsep}{2pt}
\item \textbf{\texttt{exp\_mcvc\_abilene.py}.} Data: the ABILENE topology and feature/label instances are loaded from
\end{itemize}
  \texttt{baselines/cpaior23\_branch\_and\_learn/MCVC/} (12 nodes, 15 edges, from file). Solver: ours, a batched exact
  brute-force min-cost vertex cover over all 2\textasciicircum{}12 subsets, validated against Gurobi ILP and the authors' \texttt{test.cpp}
  regret; the realized objective and penalty are ported from \texttt{test.cpp}. Classification: adapted (data and objective
  reused, exact solver reimplemented).

\begin{itemize}\setlength{\itemsep}{2pt}
\item \textbf{\texttt{interpretable\_predictor.py}.} Data: PyEPO problems (knapsack / shortest-path / portfolio / TSP) with PyEPO
\end{itemize}
  Gurobi solvers and \texttt{metric.regret}. The non-differentiable predictor is NOT sklearn: it is a custom batched hard
  axis-aligned tree (\texttt{HardTree}) initialized by a hand-written greedy CART fit, which PolyStep then refines on realized
  regret (a numerical demo shows zero gradient). Classification: original implementation (custom tree), PyEPO problems
  reused.

\begin{itemize}\setlength{\itemsep}{2pt}
\item \textbf{\texttt{exp\_spotree\_polystep.py}.} Data: the authors' real SPO-Tree shortest-path dataset (\texttt{non\_linear\_data\_dim4.p},
\end{itemize}
  Elmachtoub-Liang-McNellis). Solver: the authors' Gurobi shortest-path solver reused directly
  (\texttt{from decision\_problem\_solver import find\_opt\_decision}). Predictor: the authors' \texttt{SPOTree} / CART class
  (\texttt{from SPO\_tree\_greedy import SPOTree}); PolyStep refines the fixed-topology tree's split thresholds. Config is the
  ICML'20 paper default (n\_train 200, deg in \{2,10\}, max\_depth 3, min\_weights\_per\_node 20, CART pruning on 20\% val, 10
  reps). Classification: reused-as-is for the predictor, solver, and dataset, with the PolyStep refinement layer as the
  only addition.

\medskip\hrule\medskip

\section*{Named / real-feature benchmarks}

\begin{itemize}\setlength{\itemsep}{2pt}
\item \textbf{\texttt{exp\_warcraft.py}.} Data: the real Vlastelica/Pogancic Warcraft 12x12 "oneskin" release under
\end{itemize}
  \texttt{data/warcraft/warcraft\_shortest\_path\_oneskin/12x12/} (train/val/test maps, vertex weights, shortest paths). Solver:
  an ORIGINAL hand-written batched GPU Bellman-Ford with 8-connectivity and backtrack, validated against the dataset's
  stored optimal masks (not the authors' Dijkstra). The blackbox-differentiation autograd function follows Vlastelica
  ICLR 2020 but is reimplemented around the local solver. Feature model: a CNN, including \texttt{CombResNet18}, the paper's
  truncated ResNet18. Methods: two-stage, blackbox-differentiation (lam 20), PolyStep over a \texttt{LinearSubspace}.
  Classification: original-following-paper for both the solver and the CNN (dataset is authors' real data). This is the
  one named benchmark that deviates from strict solver reuse, the deviation being a validated reimplementation.

\begin{itemize}\setlength{\itemsep}{2pt}
\item \textbf{\texttt{exp\_tsp\_embed.py}.} Data: the real Globe-TSP release under \texttt{data/globe\_tsp/\{k\}\_countries\_from\_100/} (country
\end{itemize}
  flags, GPS, geodesic distance matrices, optimal-tour labels). Solver: REUSED-AS-IS, \texttt{gurobi\_tsp} and \texttt{\_subtour}
  copied verbatim from \texttt{martius-lab/blackbox-backprop} (\texttt{travelling\_salesman.py}), a Gurobi MIP with lazy
  subtour-elimination cuts; the authors' \texttt{TspSolver} DBB autograd function is reused with one documented fix
  (\texttt{ctx.suggested\_tours} set in forward to remove an upstream AssertionError) and a static-method API update. Feature
  model: a shared per-flag CNN predicting a 3-D coordinate inducing a Euclidean distance matrix. Methods: two-stage,
  blackbox-diff (lam 20), PolyStep. Config: k in \{5,10\}, n\_train 10000. Classification: solver reused-as-is (with a
  documented bug-fix), CNN original-following-paper.

\begin{itemize}\setlength{\itemsep}{2pt}
\item \textbf{\texttt{exp\_matching\_embed.py}.} Data: SYNTHETIC, regenerated from the paper's defined process because the original
\end{itemize}
  Edmond dataset URL is dead and the migrated file lacks the per-digit values needed for the cost metric; built from a
  local MNIST cache with the paper's two-digit edge cost \texttt{10*d\_lo + d\_hi}, the 4x4 grid edge list taken verbatim from
  the blackbox-differentiation \texttt{data/utils.py}. Solver: REUSED-AS-IS, the vendored Blossom V min-cost perfect matching
  from \texttt{blossom\_python} (built under \texttt{data/mnist\_matching/blossom\_build/}), with a networkx exact fallback; the authors'
  \texttt{MinCostPerfectMatchingSolver} DBB autograd function is reused. Feature model: a GroupNorm CNN. Classification: solver
  reused-as-is, CNN original-following-paper, data original-following-paper (forced regeneration, documented).

\begin{itemize}\setlength{\itemsep}{2pt}
\item \textbf{\texttt{exp\_amazon\_polystep.py}.} Data: the official Amazon ALMRRC 2021 challenge data, expected under
\end{itemize}
  \texttt{data/amazon\_lmrrc/} but NOT currently on disk (only \texttt{exp\_results/amazon\_lmrrc.json} from an earlier run remains, so
  the dataset must be re-fetched for reproduction). Solver: REUSED-AS-IS, the AWS reference router
  \texttt{aro.model.zone\_utils.zone\_based\_tsp} (PPM zone ordering + per-zone OR-tools TSP) from
  \texttt{baselines/amazon-routing-challenge-sol/}, scored verbatim by the official \texttt{baselines/rc-cli/scoring/score.py}.
  Feature model: a small \texttt{ZoneNet} MLP replacing the hand-crafted PPM affinity. Methods: PolyStep only as learner
  (minimizes the realized official score, label-free), with the reference PPM as baseline. Classification: solver and
  scorer reused-as-is, predictor original by design, data prep adapted (faithful node-order and travel-time rebuild).

\begin{itemize}\setlength{\itemsep}{2pt}
\item \textbf{\texttt{exp\_icon\_polystep.py}.} Data: the real ICON energy-scheduling benchmark (Irish SEMO 2013 prices \texttt{prices2013.dat},
\end{itemize}
  scheduling instances \texttt{load\{1,2,3\}/day01.txt}) via \texttt{baselines/predopt-benchmarks/Energy/Trainer}, split 550/100/142
  days (paper default). Solver: REUSED-AS-IS, \texttt{Trainer.comb\_solver.SolveICON} (exact Gurobi ILP) plus \texttt{batch\_solve},
  \texttt{regret\_fn}, \texttt{data\_reading}. Feature model: the paper's \texttt{nn.Linear(nfeat,1)}. Methods: the authors' Lightning modules
  reused verbatim (\texttt{baseline\_mse}, \texttt{SPO}, \texttt{DBB}) at their config.json LRs, plus PolyStep as the sole addition. qpth /
  cvxpy are stubbed only at import time for unused DCOL/QPTL models. Classification: reused-as-is for data, solver,
  metric, and all baselines; PolyStep added. Paper-default config.

\begin{itemize}\setlength{\itemsep}{2pt}
\item \textbf{\texttt{exp\_ieee\_polystep.py}.} Data: the real IEEE-CIS 3rd Technical Challenge phase-1 data (instances, scenarios, and
\end{itemize}
  Oct-2020 actuals \texttt{All\_data.csv}) under \texttt{baselines/mare\_optimiser/COMPETITION DATASET FILES/}. Solver: REUSED-AS-IS,
  the co-author MA\&RE Gurobi MILP optimiser (\texttt{Optimizer.py}, \texttt{Instance.py}, \texttt{Data.py}, \texttt{Setting}, \texttt{Util}, \texttt{Solution})
  from \texttt{baselines/mare\_optimiser/codes/python}. A local \texttt{schedule\_cost} helper re-implements the exact quadratic
  peak-demand charge to mirror \texttt{Optimizer.create\_objective()} for the fixed-solution evaluation path. Feature model: a
  6-parameter calendar-basis forecast-rescaling predictor (zero-init equals two-stage). Methods: two-stage and PolyStep
  (SPO+ not applicable, prediction sits in the forecast). Classification: reused-as-is for solver, data loader, and cost
  model; predictor and PolyStep loop original by design. Paper-default phase-1 instance.

\medskip\hrule\medskip

\section*{Districting (DistrictNet)}

Upstream: \texttt{https://github.com/cheikh025/DistrictNet} at HEAD \texttt{7a624a3}. The tree was compared against a fresh clone to
extract the minimal real changes. The work splits into (i) a small set of edits to tracked upstream files, (ii) a new
PolyStep estimator and train-at-scale driver, and (iii) a separate self-contained Python reimplementation that reuses
only the DistrictNet data.

\textbf{Tracked upstream-file edits (minimal, with reasons).}
\begin{itemize}\setlength{\itemsep}{2pt}
\item \texttt{experiment\_evaluator.jl}: adds \texttt{"polyStep"} to \texttt{solution\_types} and raises evaluation \texttt{NB\_SCENARIO} from 10 to 50.
\item \texttt{experiments.jl}: includes and routes the new \texttt{polyStep} module through \texttt{find\_solution} / \texttt{find\_solution\_small}.
\item \texttt{src/Estimators/\{AvgTsp,BD,FIG\}.jl}: solver \texttt{MAX\_TIME} reduced from 1200 s to 120 s (compute-budget change).
\item \texttt{src/Estimators/predictGnn.jl}: \texttt{NB\_EPOCHS} 1000 to 40, \texttt{MAX\_TIME} 1200 to 120.
\item \texttt{src/districtNet.jl}: \texttt{NB\_EPOCHS} 100 to 40 (env \texttt{DN\_EPOCHS}), \texttt{MAX\_TIME} 1200 to 120; adds a \texttt{DN\_UNTRAINED} random-GNN
\end{itemize}
  ablation branch and \texttt{[TIMING]} instrumentation accumulating build / load / exact-solve / feature times.
\begin{itemize}\setlength{\itemsep}{2pt}
\item \texttt{src/learning.jl}: per-epoch timing plus CMST set-partition solve counters (\texttt{SPO\_CALLS}, \texttt{SPO\_TIME}).
\item \texttt{src/Solver/exactsolver.jl}: wraps \texttt{set\_partitioning\_optimizer} in a try/finally that increments \texttt{SPO\_CALLS} and
\end{itemize}
  accumulates \texttt{SPO\_TIME} (timing only, logic unchanged).
\begin{itemize}\setlength{\itemsep}{2pt}
\item \texttt{src/district.jl}: two robustness guards. (1) \texttt{compute\_cost\_via\_SAA} retries the C++ evaluator up to 6 times with
\end{itemize}
  backoff, because concurrent scenario-file regeneration across parallel jobs can let a read catch a mid-write JSON file.
  (2) \texttt{compute\_cost\_via\_CMST} adds a disconnected-district guard: if \texttt{prim\_mst} returns fewer than n-1 edges the blocks
  are not mutually reachable, so the district is penalized like a singleton instead of summing over an empty edge set
  (which would throw).
\begin{itemize}\setlength{\itemsep}{2pt}
\item \texttt{src/Solver/Kruskal.jl}: +249 lines implementing \textbf{cascading chain-repair} for size feasibility. \texttt{repair\_districts}
\end{itemize}
  now escalates: plain grow/shrink, then constraint-aware cascade repair on both sides (grow an undersized orphan, or
  shed an over-max district, by relocating one node along a BFS path of adjacent districts to or from a district with
  spare capacity, keeping the FULL district count even under the 2*min > max no-split window), then a last-resort
  \texttt{dissolve\_min\_districts} only if no reachable donor/receiver exists. A \texttt{DN\_NO\_REPAIR} env exposes the raw pre-repair
  orphan for the before/after figures. This is the change that makes large instances (e.g. ilsdefrance) feasible.

\textbf{New PolyStep estimator: \texttt{src/Estimators/polyStep.jl} (untracked, 375 lines).} A gradient-free,
softmax-barycentric (optimal-transport-style) trainer that minimizes the REALIZED districting cost directly, with no
imitation targets and no Fenchel-Young loss. It reuses the DistrictNet GNN architecture and forward pass
(\texttt{STRATEGY = "districtNet"}, \texttt{build\_gnn\_model}, \texttt{predict\_theta}) and the existing CMST forward solver. Two variants:
\texttt{GNNtrainer\_polyStep} (small 30-BU / target-3 instances, scored with the cached district costs, no C++ at train time)
and \texttt{GNNtrainer\_polyStep\_atscale} (deploy-scale real cities, scored with the true C++ SAA evaluator and a memoized
district-cost cache, with no labels and no exact set-partition at any point). An \texttt{PS\_SIZE\_PENALTY} env adds an explicit
district-size penalty straight into the scalar objective (the An \#8 feasibility-aware variant), which the
differentiable PerturbedAdditive(CMST) path cannot carry because PolyStep consumes only the scalar cost. The estimator
also dumps a per-step replay JSON for visualization. Classification: original PolyStep estimator following the
project's own method, reusing DistrictNet's GNN and CMST machinery.

\textbf{New driver: \texttt{trainatscale\_driver.jl} (untracked, 159 lines).} Compares three deploy pipelines on held-out large
real cities (Marseille / Birmingham fully held out, London / Leeds scaled up): DistrictNet trained small (FY label
imitation), PolyStep trained small (realized cost), and PolyStep trained at scale (target-20, the new variant). Metric
is the shared C++ SAA routing cost plus a feasibility check (full district count, no size violations), with each
solution evaluated by the same module that produced it. Plus assorted \texttt{run\_dn\_*.sh} and \texttt{*\_driver.jl} shell/Julia
drivers (city redeploy, "does data help" untrained-vs-trained, penalty training, raw-orphan capture, side-by-side
feasibility, timing).

\textbf{Python launcher: \texttt{districtnet\_polystep.py}.} A SEPARATE self-contained reimplementation that reuses only the
DistrictNet data and its precomputed realized-cost cache (no Julia / C++ build): each city's \textasciitilde{}450 candidate districts
carry a cached expected routing cost, a set-partition ILP is solved on predicted costs with Gurobi, and the TRUE cached
cost is incurred. Trains a geometric-feature predictor with two-stage, SFGE, and PolyStep (pristine import).
Classification: original implementation reusing the authors' data and cache, with a Gurobi set-partition solver.

\textbf{Data.} The authors' real DistrictNet data is present: 211 geojson city blocks under \texttt{data/geojson/}, precomputed
\texttt{data/tspCosts/} cost files, \texttt{data/cities}, \texttt{data/maps}, \texttt{data/maptiles}.

\textbf{Paper-default caveat.} The DistrictNet runs deviate from the paper-default compute budget: solver / ILS \texttt{MAX\_TIME}
was cut from 1200 s to 120 s across estimators, GNN epochs were reduced (predictGnn 1000 to 40, districtNet 100 to 40),
and evaluation \texttt{NB\_SCENARIO} was set to 50 (the upstream evaluator default was 10). The model architectures, the FY
training objective, the CMST forward solver, and the data are otherwise the authors' originals.

\medskip\hrule\medskip

\section*{Inventory routing (IRP)}

Two distinct IRP threads exist: the "Toni" column from the authors' Julia DSIRP code, and the project's own Python
predict to CPCTSP to simulate pipeline.

\textbf{Toni / DSIRP authors' code} at \texttt{/media/anindex/Data/project-cache/ot-or-project/InferOpt\_DSIRP} (upstream
\texttt{tonigreif/InferOpt\_DSIRP}). This is a copied cache, not a git checkout (no \texttt{.git}), so changes were located by reading
the source directly.
\begin{itemize}\setlength{\itemsep}{2pt}
\item \texttt{src/sirp\_model.jl}: a \textbf{uniform-demand branch} was added. The uniform-10 instances ship no \texttt{mean\_demand}, so for
\end{itemize}
  \texttt{pattern == "uniform"} the start inventory is set to \texttt{0.75 * max\_inventory} (lines 181-184) and the sampling demand is
  \texttt{Uniform(0, 0.5 * max\_inventory)} (lines 218-219), matching the authors' stated demand process. The normal, bimodal,
  and contextual branches are unchanged.
\begin{itemize}\setlength{\itemsep}{2pt}
\item \texttt{src/auxiliar.jl}: the current \texttt{grb\_model} uses a single shared persistent \texttt{GRB\_ENV} and sets only \texttt{OutputFlag = 0}.
\end{itemize}
  No explicit Gurobi \texttt{Threads} attribute is present in the working copy; the previously noted Threads 1 to 8 change is
  not observable here (it appears to have been reverted or never landed in this copy). This is reported as found.
\begin{itemize}\setlength{\itemsep}{2pt}
\item Training uses the authors' pipeline at paper defaults: \texttt{train\_pipeline.jl} defaults are \texttt{paradigm = dagger},
\end{itemize}
  \texttt{pipeline\_epochs = 150}, \texttt{dagger\_iterations = 20}, \texttt{look\_ahead = 6}, \texttt{shortage\_penalty = 200}, \texttt{milp\_solver = gurobi},
  with the 12-quantile statistical model and FYL defaults (5 samples, epsilon 20). The Toni column consumed by the
  Python head-to-head is the authors' trained GLM weights, loaded unchanged into a dimension-matched predictor.

\textbf{Project IRP scripts} (following Greif/Bouvier/Parmentier, arXiv:2402.04463).
\begin{itemize}\setlength{\itemsep}{2pt}
\item \textbf{\texttt{exp\_irp\_polystep.py}.} Data: the authors' RELEASED instances loaded from JSON (the DSIRP \texttt{instances/} directory),
\end{itemize}
  parsed against the authors' schema. Solver: an ORIGINAL Python reimplementation, \texttt{cpctsp} is our own Gurobi MILP
  mirroring \texttt{src/pctsp.jl}, the PINN predictor mirrors \texttt{src/stat\_model.jl}, and the rollout mirrors
  \texttt{src/evaluation\_pipeline.jl}. Methods: a GLM-init PINN baseline (authors' init \texttt{w\_inv = -1/72}, \texttt{w\_pen = +1/72}) and
  PolyStep (pristine, orthoplex, num\_probe 1). Paper-default constants carried explicitly: 12 quantiles, look-ahead 6,
  shortage penalty 200, eval horizon 10. Classification: original-following-paper pipeline plus pristine PolyStep.
\begin{itemize}\setlength{\itemsep}{2pt}
\item \textbf{\texttt{irp\_gpu.py}.} An ORIGINAL GPU-batched EXACT CPCTSP solver: exact Held-Karp DP precomputes the optimal tour length
\end{itemize}
  for every customer subset once per instance, then each CPCTSP is a batched membership-matmul + capacity-mask + argmax
  over all subsets; exactness is self-verified against the gurobipy reference (\texttt{verify\_vs\_gurobi}). This is the \textasciitilde{}47x
  speedup component. Classification: original, validated equal-objective to Gurobi.
\begin{itemize}\setlength{\itemsep}{2pt}
\item \textbf{\texttt{exp\_irp\_headtohead.py}.} Reuses the authors' released instances and the authors' trained Toni GLM weights (loaded
\end{itemize}
  from the DSIRP cache \texttt{training/solutions/dagger/.../penalty\_200}), with \texttt{look\_ahead} forced to match the value the
  weights were trained at. Four columns: untrained PINN, Toni-FY (authors' weights in our identical simulator), PolyStep
  (pristine, GPU-batched closure), and SFGE (our score-function trainer over the rollout). Solvers: our Gurobi CPCTSP
  for the eval path and our GPU-batched exact CPCTSP for training. Classification: original-following-paper harness that
  reuses the authors' instances and trained weights as-is.
\begin{itemize}\setlength{\itemsep}{2pt}
\item \textbf{\texttt{\_diag\_irp.py}.} A diagnostic over one fixed authors' instance and one Toni solutions JSON, printing untrained
\end{itemize}
  versus Toni-FY held-out cost. Reuses the \texttt{exp\_irp\_polystep} scaffold; no training, no PolyStep. Classification:
  original diagnostic consuming authors' instance and weights as-is.

\medskip\hrule\medskip

\section*{Oracle-cost theory and scaling}

These are original measurement and theory harnesses over the shared \texttt{pto/} problems (PyEPO data + PyEPO baselines +
our GPU solvers + pristine PolyStep). No external authors' experiment code is reused.

\begin{itemize}\setlength{\itemsep}{2pt}
\item \textbf{\texttt{exp\_oracle\_crossover.py}.} Measures the wall-clock crossover between forward-only methods (PolyStep / SFGE,
\end{itemize}
  evaluation oracle) and SPO+ (optimization oracle), and validates the analytic cost-model quantities kappa, rho, and
  the crossover latency. A \texttt{LatencyOracle} injects a controlled busy-wait into the identical SPO+(batched) algorithm so
  the predicted crossover can be traced with byte-identical output (regret invariant). Problems: PyEPO TSP and knapsack.
  SPO+ reused from \texttt{exp\_fair\_batched\_spo} (PyEPO algorithm). Classification: original cost-model theory harness.

\begin{itemize}\setlength{\itemsep}{2pt}
\item \textbf{\texttt{exp\_scaling.py}.} Wall-clock and regret versus problem scale on PyEPO TSP and knapsack, demonstrating the
\end{itemize}
  per-instance-exact-solve (SPO+) versus batched-forward (PolyStep / SFGE) asymmetry, with identical two-stage warm
  start. SPO+ via the PyEPO wrapper. Classification: original scaling harness.

\begin{itemize}\setlength{\itemsep}{2pt}
\item \textbf{\texttt{exp\_scaling\_memory.py}.} Peak GPU memory and per-step wall-clock versus predictor parameter count, finding where
\end{itemize}
  SFGE (backprop activation graph) OOMs but PolyStep (forward-only, chunkable) survives. Uses its own minimal synthetic
  top-k task rather than \texttt{pto/problems.py}. The real \texttt{PolyStepOptimizer} (pristine) and a real SFGE step are exercised,
  with single-chunk peak estimates validated against an actual \texttt{pso.step}. Classification: original memory-measurement
  harness (MeZO-style).

\begin{itemize}\setlength{\itemsep}{2pt}
\item \textbf{\texttt{exp\_config\_sweep.py}.} Normalized regret and wall-clock across the PolyStep config grid (polytope x num\_probe x
\end{itemize}
  seeds) over both PyEPO problems and the constraint-prediction families (\texttt{MDKPConsumption}, \texttt{KnapsackWeights}), to
  derive the recommended default polytope and probe count. PolyStep only, pristine, warm-started from two-stage.
  Classification: original config-robustness harness.

\medskip\hrule\medskip

\section*{Compact per-experiment table}

{\footnotesize\begin{longtable}{p{0.16\textwidth}p{0.18\textwidth}p{0.13\textwidth}p{0.30\textwidth}p{0.11\textwidth}}
\toprule
\textbf{Experiment} & \textbf{Upstream source} & \textbf{Reuse class} & \textbf{Key local changes} & \textbf{Paper-default config?} \\ \midrule \endhead
challenge\_established.py & PyEPO (data, models, SPO+/PFYL/IMLE) & reused-as-is + original SFGE & shared two-stage warm start; PolyStep pristine & Yes (PyEPO defaults) \\
exp\_fair\_batched\_spo.py & PyEPO SPO+ & reused-as-is + adapted backend & batched SPO+ twin mirroring SPOPlusFunc, checked vs Gurobi & Yes \\
exp\_polystep\_vs\_sfge.py & PyEPO SP/knapsack & reused-as-is + original SFGE & per-method tuned HPs from prior sweep & Yes (PyEPO defaults) \\
exp\_polystep\_bench.py & PyEPO four problems + full DFL suite & reused-as-is & PolyStep HPs from tune JSON; Wilcoxon & Yes \\
exp\_polystep\_tune.py & PyEPO setups & reused-as-is (PolyStep utility) & tuning grid only & N/A (tuning) \\
exp\_polystep\_ablation.py & PyEPO setups & reused-as-is (PolyStep utility) & polytope x probe sweep & N/A (ablation) \\
exp\_sfge\_suite.py & SFGE paper (Silvestri JAIR'24) & original-following-paper & own solvers (frac-knapsack, set-multicover) vs cvxpy/Gurobi & Follows paper setup \\
exp\_spo\_convergence.py & PyEPO SPO+ & reused-as-is + original SFGE & epoch-grid convergence sweep & Yes \\
cross\_sfge.py & PyEPO problems + SPO+ & reused-as-is + original SFGE & all warm-started from two-stage & Yes \\
fair\_init.py & PyEPO knapsack + SPO+ & reused-as-is (control) & cold-vs-warm init integrity check & Yes \\
blackbox\_constraint.py & PyEPO knapsack generator & original-following-paper & prediction-in-constraint variant; our DP solver & PyEPO defaults \\
exp\_constraint\_complexity.py & none (own synthetic) & original & L0-L4 ladder; greedy solver Gurobi-verified & N/A (original) \\
exp\_constraint\_scaling.py & own MDKP synth + PyEPO contrast & original + reused-as-is & MDKPConsumption regime + PyEPO SPO+ contrast & PyEPO contrast at defaults \\
exp\_manyconstraints.py & PyEPO SP + own MDKP & adapted + original & two regimes; solve/Gurobi accounting & PyEPO SP at defaults \\
exp\_odece\_full.py & ODECE generators + Results/ & adapted & ODECE generator reused; ODECE solver/loss NOT run; numbers cited & Yes (ODECE scripts) \\
exp\_odece\_mdkp.py & ODECE genCapacity & adapted & ODECE generator reused; our greedy solver + Gurobi oracle & Yes (MDKP\_CapaExp.sh) \\
exp\_maxflow\_polska.py & cpaior23 B\&L POLSKA files & adapted & topology from file; corrA ported; batched solver ours & Follows B\&L setup \\
exp\_mcvc\_abilene.py & cpaior23 B\&L ABILENE files & adapted & data + objective from file; exact 2\textasciicircum{}12 solver ours & Follows B\&L setup \\
interpretable\_predictor.py & PyEPO problems & original-impl & custom HardTree + CART refined by PolyStep (not sklearn) & PyEPO defaults \\
exp\_spotree\_polystep.py & SPOTree (ICML'20) code + data & reused-as-is & authors' SPOTree/CART + Gurobi SP; PolyStep refines thresholds & Yes (ICML'20 default) \\
exp\_warcraft.py & Vlastelica Warcraft data & original-following-paper & dataset authors'; solver+CNN reimplemented, validated vs optima & Yes (12x12 defaults) \\
exp\_tsp\_embed.py & blackbox-backprop TSP solver + Globe-TSP data & reused-as-is & vendored gurobi\_tsp/TspSolver; one documented bug-fix & Yes \\
exp\_matching\_embed.py & Blossom V + MNIST process & reused-as-is solver, original data & Blossom vendored; dataset regenerated (dead URL) & Follows paper cost \\
exp\_amazon\_polystep.py & AWS ALMRRC router + rc-cli scorer & reused-as-is & OR-tools router + official scorer verbatim; ZoneNet MLP added; data absent on disk & Yes (official protocol) \\
exp\_icon\_polystep.py & predopt-benchmarks ICON & reused-as-is & SolveICON + all baselines verbatim; PolyStep added & Yes (550/100/142) \\
exp\_ieee\_polystep.py & MA\&RE optimiser + IEEE-CIS data & reused-as-is & Gurobi MILP + loaders verbatim; calendar predictor + PolyStep & Yes (phase-1) \\
districtnet\_polystep.py & DistrictNet data + cache & original-impl reusing data & Gurobi set-partition ILP; two-stage/SFGE/PolyStep on geometric features & Reuses authors' data \\
DistrictNet/*.jl drivers & cheikh025/DistrictNet & adapted & new polyStep.jl estimator; cascade chain-repair; train-at-scale driver; timing & Budgets reduced (see caveat) \\
exp\_irp\_headtohead.py & DSIRP instances + Toni weights & original-following-paper harness & authors' instances + trained weights reused; our Gurobi/GPU CPCTSP & Yes (look-ahead, penalty 200) \\
exp\_irp\_polystep.py & DSIRP instances & original-following-paper & Python port of predict->CPCTSP->simulate; pristine PolyStep & Yes (12 quantiles, LA 6) \\
irp\_gpu.py & none (own solver) & original & GPU-batched exact CPCTSP, validated vs Gurobi & N/A (accelerator) \\
\_diag\_irp.py & DSIRP instance + Toni weights & original diagnostic & sanity-check of loaded weights & Inherits IRP constants \\
exp\_oracle\_crossover.py & pto harness (PyEPO) & original & cost-model theory + latency oracle & N/A (own theory) \\
exp\_scaling.py & pto harness (PyEPO) & original & wall-clock/regret vs scale & N/A \\
exp\_scaling\_memory.py & none (own task) & original & memory crossover; own top-k task & N/A \\
exp\_config\_sweep.py & pto harness (PyEPO + own MDKP) & original & polytope/probe robustness sweep & N/A (derives default) \\
\bottomrule\end{longtable}}

\medskip\hrule\medskip

\section*{Notes and items flagged}

\begin{itemize}\setlength{\itemsep}{2pt}
\item \textbf{DSIRP Gurobi Threads change not observed.} The brief named a Threads 1 to 8 edit in
\end{itemize}
  \texttt{InferOpt\_DSIRP/src/auxiliar.jl}; the current working copy of \texttt{grb\_model} sets only \texttt{OutputFlag = 0} on a shared
  persistent environment, with no Threads attribute anywhere in the repo. Only the uniform-demand branch in
  \texttt{sirp\_model.jl} is confirmed as a local change. Reported as found.

\begin{itemize}\setlength{\itemsep}{2pt}
\item \textbf{DistrictNet compute budgets} were reduced below paper defaults (MAX\_TIME 1200 to 120 s, epochs cut, eval
\end{itemize}
  NB\_SCENARIO 10 to 50). The architectures, FY objective, CMST solver, and data are the authors' originals. This is the
  one place where the runs depart from the strict paper-default-config requirement, and it is a budget choice rather
  than a method change.

\begin{itemize}\setlength{\itemsep}{2pt}
\item \textbf{Two named benchmarks depart from strict author-solver reuse, both documented and validated:} Warcraft (solver and
\end{itemize}
  CNN reimplemented from the paper, validated against the authors' stored optimal masks) and MNIST matching (dataset
  regenerated from the paper's process because the authors' file is dead and label-incomplete, with the Blossom V solver
  still vendored).

\begin{itemize}\setlength{\itemsep}{2pt}
\item \textbf{One dataset is currently missing from disk:} the Amazon ALMRRC challenge data expected at \texttt{data/amazon\_lmrrc/} is
\end{itemize}
  absent (only the earlier result JSON remains), so that experiment needs the dataset re-fetched to reproduce. The
  solver, scorer, and code path are all present and reused verbatim.

\begin{itemize}\setlength{\itemsep}{2pt}
\item Every experiment was classifiable; none was left unclassified.
\end{itemize}

\end{document}